% Generated by src/equivalent_rights_proposals.jl -- do not edit by hand.
\begin{table}[htbp]
\centering
\small
\caption{Club members left below their welfare under the proposal, by variant}
\begin{tabular}{lcccc}
  \toprule
  & \multicolumn{2}{c}{\textbf{Wolfram et al.}} (of 47) & \multicolumn{2}{c}{\textbf{Banerjee et al.}} (of 179) \\
  \cmidrule(lr){2-3} \cmidrule(lr){4-5}
  \textbf{Variant} & \textbf{isolated} & \textbf{joint} & \textbf{isolated} & \textbf{joint} \\
  \midrule
  \multicolumn{5}{l}{\textit{Members losing on equally-distributed-equivalent consumption}} \\
  \quad Variant 1 & 24 & 1 & 38 & 33 \\
  \quad Variant 2 & 2 & 2 & 37 & 37 \\
  \quad Variant 3 & 1 & 2 & 41 & 33 \\
  \quad Variant 4 & 6 & 9 & 51 & 70 \\
  \midrule
  \multicolumn{5}{l}{\textit{Members losing on mean consumption per capita}} \\
  \quad Variant 1 & 29 & 0 & 37 & 0 \\
  \quad Variant 2 & 0 & 0 & 0 & 0 \\
  \quad Variant 3 & 0 & 3 & 38 & 24 \\
  \quad Variant 4 & 2 & 7 & 10 & 43 \\
  \bottomrule
\end{tabular}
\label{tab:equiv_rights_losers}
\\[4pt]
{\footnotesize Note: a member counts as losing when its NPV of consumption over 2030--2100 falls short of what it obtains under the proposal itself. Variant 1 makes every member indifferent by construction, so any count there is numerical noise. Variants 2 to 4 return the same surplus of rights under different sharing rules: a split chosen to minimise the largest relative shortfall (2), one proportional to population times marginal utility (3), and uniform scaling (4).}
\end{table}
